\section{Segment Intersection}

Given two line segments $\overline{AB}$ and $\overline{CD}$, we want to
determine whether they intersect. This operation is fundamental to collision
detection, polygon validation, visibility testing, and motion planning.

\subsection{Proper Intersections}

A proper intersection occurs when two line segments cross through the
interiors of both segments.

Let the two segments be

\[
\overline{AB}
\qquad\text{and}\qquad
\overline{CD}.
\]

Define

\[
o_1 = \operatorname{orientation}(A,B,C),
\]

\[
o_2 = \operatorname{orientation}(A,B,D),
\]

\[
o_3 = \operatorname{orientation}(C,D,A),
\]

and

\[
o_4 = \operatorname{orientation}(C,D,B).
\]

The segments properly intersect when the endpoints of each segment lie on
opposite sides of the other segment:

\[
o_1 o_2 < 0
\]

and

\[
o_3 o_4 < 0.
\]

A proper intersection therefore excludes endpoint touching and collinear
overlap. These cases may still count as general segment intersections, but
they are geometrically different from a crossing through both segment
interiors.

This distinction is useful for visibility testing. A candidate visibility
edge should be rejected when it properly crosses an obstacle boundary, while
some forms of boundary touching must remain allowed.

\subsection{Degenerate Cases}

The proper-intersection test alone does not handle endpoint touching,
collinear overlap, or one endpoint lying on another segment. These cases
require a point-on-segment test.

A point $P$ lies on $\overline{AB}$ if
\[
\operatorname{orientation}(A,B,P)=0
\]
and
\begin{align*}
\min(x_A,x_B)&\leq x_P\leq\max(x_A,x_B),\\
\min(y_A,y_B)&\leq y_P\leq\max(y_A,y_B).
\end{align*}

In the current project convention, endpoint touching and collinear overlap
both count as intersection.

\subsection{Complexity}

The segment-intersection test performs a constant number of arithmetic
operations. Therefore,
\[
T(n)=O(1).
\]

Although simple, this predicate becomes one of the most frequently used
operations in later motion-planning algorithms.
